\section{Levantamento de Requisitos}

Em coerência com a Engenharia de Requisitos discutida na Base
Teórica, o levantamento apresentado nesta seção corresponde
às etapas de \textit{Análise} do ciclo de BPM: a partir do
diagnóstico do processo atual, busca-se identificar os
gargalos operacionais e transformá-los em requisitos capazes
de orientar o redesenho do processo, apresentado na seção
seguinte. O levantamento de requisitos foi realizado a partir
das respostas do formulário estruturado aplicado aos três
participantes definidos na metodologia: uma Analista de
Qualidade e Clientes, uma Supervisora de Qualidade e um
Gerente da Área. Os três participantes registraram
concordância voluntária com a pesquisa. As perguntas foram
segmentadas conforme o nível de atuação, de modo que os
campos sem resposta nos blocos específicos não representam
omissão, mas sim questões direcionadas a outro perfil
organizacional.

A análise das respostas foi conduzida por agrupamento
temático, considerando cinco eixos: processo atual,
dificuldades e riscos, regras de negócio, funcionalidades
esperadas e benefícios gerenciais. A partir desses eixos,
as necessidades relatadas foram convertidas em requisitos
funcionais, requisitos não funcionais e regras de negócio.
Para a priorização foi utilizada a técnica MoSCoW,
considerando como \textit{Must Have} os requisitos
necessários para que o fluxo opere com rastreabilidade e
governança mínimas, como \textit{Should Have} aqueles que
aumentam a qualidade operacional sem impedir a operação
básica, e como \textit{Could Have} as funcionalidades de
evolução que podem ser incorporadas em ciclos posteriores.

\subsection{Diagnóstico do Processo Atual}

A resposta da Analista de Qualidade e Clientes permitiu
caracterizar o processo operacional utilizado antes da
solução proposta. Os chamados internos eram registrados no
SAP Service, ferramenta também utilizada como CRM para
atendimentos e reclamações de clientes externos. Para
viabilizar esse registro, as ocorrências internas eram
associadas a um cadastro denominado ``Experiência do
Cliente'', no qual era criado um ticket e realizado o
direcionamento à área responsável. Paralelamente ao registro
no sistema, a chegada do chamado era comunicada manualmente
aos responsáveis por meio de um grupo de WhatsApp.

Esse modelo apresentava duas fragilidades centrais. A
primeira estava relacionada à mistura, em um mesmo ambiente,
de chamados internos e manifestações externas de clientes,
o que poderia distorcer a leitura de volumes e dificultar a
governança das ocorrências internas. A segunda estava
relacionada à dependência de atividades manuais para
comunicação e acompanhamento, especialmente quanto à
obtenção de retorno das áreas responsáveis.

A Analista também diferenciou os dois tipos de chamado
utilizados no processo. O chamado interno corresponde a uma
problemática organizacional que necessita de análise e
tratativa por outra área. O chamado massivo corresponde a
uma ocorrência de mesmo problema capaz de impactar três ou
mais clientes, normalmente associada a falhas técnicas ou
sistêmicas com maior potencial de impacto. Essa distinção
passou a ser considerada na modelagem do processo e nas
regras de priorização.

Na perspectiva da Supervisora de Qualidade, a proposta de
um ambiente específico para registro e acompanhamento foi
considerada aderente às necessidades operacionais. A
respondente destacou como benefícios esperados a maior
agilidade no direcionamento, a rastreabilidade das
tratativas, a simplificação do processo em comparação ao
SAP Service, a notificação imediata da área responsável e a
possibilidade de consolidação das informações para análise
gerencial. Como pontos de atenção, foram identificadas a
necessidade de registrar atualizações ou justificativas
quando a solução depender de fatores externos e a
necessidade de permitir a transferência do chamado quando o
responsável inicialmente definido não for o correto.

Na perspectiva gerencial, foram confirmadas como principais
fricções a dificuldade de governança dos chamados internos,
a falta de padronização dos registros e do direcionamento e
a dificuldade das áreas executoras em visualizar suas
pendências. O Gerente da Área avaliou que o fluxo proposto
possui aderência aos objetivos estratégicos da gerência e
da organização, destacando como benefícios esperados a
rastreabilidade, a padronização, a disponibilidade de
indicadores gerenciais e a redução do tempo de resolução.
O fluxo foi aprovado, em princípio, sem ressalvas.

A Tabela~\ref{tab:sintese-levantamento} sintetiza as
contribuições de cada perfil para a especificação da
solução.

\begin{table}[H]
\caption{Síntese do levantamento de requisitos por perfil}
\label{tab:sintese-levantamento}
\centering
\begin{tabular}{p{2.4cm} p{5cm} p{5cm}}
\hline
\textbf{Perfil} & \textbf{Principais evidências} &
\textbf{Implicações para a solução} \\
\hline
Analista de Qualidade e Clientes &
Uso do SAP Service para ocorrências internas e externas;
comunicação complementar por WhatsApp; dificuldade de
retorno; necessidade de protocolo, evidências, regras,
causa raiz, solução, SLA e indicadores. &
Criar canal próprio; padronizar a abertura; gerar protocolo;
automatizar notificações; registrar causa raiz e ação;
separar chamados internos de manifestações externas. \\
\hline
Supervisora de Qualidade &
Necessidade de classificação e solução por prazo;
rastreabilidade; transferência de responsável; registro de
atualizações e justificativas quando o SLA não puder ser
cumprido. &
Definir papéis de classificação; controlar mudança de
responsável; manter histórico; prever mecanismo de
justificativa e acompanhamento; estruturar regras de SLA. \\
\hline
Gerente da Área &
Dificuldade de governança, padronização e acesso às
pendências; necessidade de indicadores; aprovação do fluxo
proposto. &
Centralizar os registros; disponibilizar visão gerencial;
priorizar volume, tempo médio de resolução e distribuição
por área; garantir aderência ao processo organizacional. \\
\hline
\end{tabular}
\raggedright
\par
{\fontsize{10}{12}\selectfont Fonte: Elaborado pelo autor
com base na pesquisa de campo}
\end{table}

\subsection{Requisitos Funcionais}

Os requisitos funcionais representam os serviços que a
solução deve disponibilizar aos usuários e às áreas
envolvidas no processo. A Tabela~\ref{tab:requisitos-funcionais}
apresenta os requisitos consolidados a partir das respostas
dos participantes.

\begin{table}[H]
\caption{Requisitos funcionais da solução}
\label{tab:requisitos-funcionais}
\centering
\begin{tabular}{p{1.0cm} p{8.7cm} p{2.3cm}}
\hline
\textbf{ID} & \textbf{Requisito funcional} &
\textbf{Prioridade} \\
\hline
RF01 & Disponibilizar um canal único e específico para
abertura de chamados internos, separado do ambiente de
reclamações externas de clientes. & Must Have \\
\hline
RF02 & Permitir que usuários autorizados de diferentes
áreas registrem novos chamados internos. & Must Have \\
\hline
RF03 & Permitir a identificação do tipo de chamado,
diferenciando, no mínimo, chamados Internos e Massivos. &
Must Have \\
\hline
RF04 & Registrar de forma padronizada as informações
necessárias à abertura, incluindo empresa, tipo, assunto,
descrição do problema e evidências, além de dados
operacionais aplicáveis ao caso. & Must Have \\
\hline
RF05 & Gerar automaticamente um número único de protocolo
para cada chamado registrado. & Must Have \\
\hline
RF06 & Confirmar a abertura ao solicitante e informar o
protocolo e o link de acesso ao registro. & Must Have \\
\hline
RF07 & Notificar a equipe de Qualidade/SAC sobre novos
chamados que aguardam classificação e direcionamento. &
Must Have \\
\hline
RF08 & Permitir à Qualidade/SAC definir prioridade, área
responsável, tipologia, responsável pela tratativa e status
do chamado. & Must Have \\
\hline
RF09 & Notificar automaticamente o responsável quando um
chamado for atribuído para sua tratativa. & Must Have \\
\hline
RF10 & Permitir a transferência do chamado para outro
responsável, preservando a rastreabilidade da alteração. &
Must Have \\
\hline
RF11 & Permitir o registro de causa raiz e ação corretiva
pela área responsável pela tratativa. & Must Have \\
\hline
RF12 & Encaminhar o chamado tratado para análise da
Qualidade/SAC antes do encerramento definitivo. &
Must Have \\
\hline
RF13 & Permitir à Qualidade/SAC concluir ou cancelar o
chamado, registrando o motivo do cancelamento quando
aplicável. & Must Have \\
\hline
RF14 & Notificar o solicitante e os responsáveis quando o
chamado for concluído ou cancelado. & Must Have \\
\hline
RF15 & Controlar os prazos de classificação e solução de
acordo com o tipo e a prioridade do chamado, possibilitando
o acompanhamento do SLA. & Must Have \\
\hline
RF16 & Permitir o registro de comentários, atualizações ou
justificativas para chamados cuja tratativa dependa de
fatores externos ou ultrapasse o prazo inicialmente
definido. & Should Have \\
\hline
RF17 & Permitir o registro da quantidade de protocolos
relacionados quando a ocorrência afetar múltiplos clientes
ou registros. & Should Have \\
\hline
RF18 & Disponibilizar categorias ou tipologias reutilizáveis
para apoiar a classificação e a análise de recorrências. &
Should Have \\
\hline
RF19 & Disponibilizar indicadores gerenciais de volume,
tempo médio de resolução, distribuição por área, tipos de
problema, chamados tratados e cumprimento de SLA. &
Must Have \\
\hline
\end{tabular}
\raggedright
\par
{\fontsize{10}{12}\selectfont Fonte: Elaborado pelo autor
com base na pesquisa de campo}
\end{table}

\subsection{Requisitos Não Funcionais}

Além das funcionalidades, o levantamento evidenciou
características de qualidade necessárias para que a solução
seja incorporada à rotina organizacional. Esses requisitos
estão apresentados na Tabela~\ref{tab:requisitos-nao-funcionais}.

\begin{table}[H]
\caption{Requisitos não funcionais da solução}
\label{tab:requisitos-nao-funcionais}
\centering
\begin{tabular}{p{1.1cm} p{8.5cm} p{2.4cm}}
\hline
\textbf{ID} & \textbf{Requisito não funcional} &
\textbf{Prioridade} \\
\hline
RNF01 & A solução deve centralizar os registros em uma fonte
única e permitir rastreabilidade desde a abertura até o
encerramento. & Must Have \\
\hline
RNF02 & A interface de abertura deve ser simples, objetiva e
adequada ao uso em ambiente corporativo, reduzindo o número
de etapas necessárias para registrar uma ocorrência. &
Must Have \\
\hline
RNF03 & A solução deve utilizar, prioritariamente, os
recursos do ecossistema Microsoft 365 já adotado pela
organização, favorecendo integração e manutenção. &
Must Have \\
\hline
RNF04 & O acesso aos registros e às ações do processo deve
respeitar as permissões corporativas e os diferentes papéis
de solicitante, Qualidade/SAC e área responsável. &
Must Have \\
\hline
RNF05 & As alterações relevantes de status, responsável,
tratativa e encerramento devem permanecer auditáveis no
histórico do registro. & Must Have \\
\hline
RNF06 & As automações devem evitar notificações duplicadas e
reprocessamentos causados pelas próprias atualizações
realizadas no item. & Must Have \\
\hline
RNF07 & A interface incorporada ao SharePoint deve adaptar-se
ao espaço disponível na página, preservando legibilidade e
usabilidade. & Should Have \\
\hline
\end{tabular}
\raggedright
\par
{\fontsize{10}{12}\selectfont Fonte: Elaborado pelo autor}
\end{table}

\subsection{Regras de Negócio}

As respostas da pesquisa permitiram definir regras
operacionais que orientam o comportamento do fluxo e a
distribuição de responsabilidades entre os participantes.
A Tabela~\ref{tab:regras-negocio} apresenta as principais
regras identificadas.

\begin{table}[H]
\caption{Regras de negócio identificadas}
\label{tab:regras-negocio}
\centering
\begin{tabular}{p{1.0cm} p{10.9cm}}
\hline
\textbf{ID} & \textbf{Regra de negócio} \\
\hline
RN01 & Todo chamado deve ser registrado assim que a
problemática for identificada, garantindo rastreabilidade
desde o início do processo. \\
\hline
RN02 & A Qualidade/SAC é responsável pela classificação e
pelo direcionamento inicial do chamado para a área
responsável pela tratativa. \\
\hline
RN03 & Chamados Massivos devem possuir tratamento mais
restritivo em razão do potencial impacto sobre maior
quantidade de clientes e indicadores. \\
\hline
RN04 & Na pesquisa de campo, a classificação proposta para
Chamados Massivos foi Alta e para Chamados Internos foi
Média. \\
\hline
RN05 & A classificação inicial de um Chamado Massivo deve
ocorrer em até 30 minutos após a abertura; para Chamados
Internos, em até um dia após a abertura. \\
\hline
RN06 & A solução de um Chamado Massivo deve ocorrer em até
24 horas após a classificação; para Chamados Internos, em
até três dias após a classificação. \\
\hline
RN07 & A área responsável deve registrar, no mínimo, a causa
raiz identificada e a ação corretiva adotada antes de
submeter a tratativa para análise da Qualidade/SAC. \\
\hline
RN08 & A conclusão definitiva do chamado deve ser realizada
pela Qualidade/SAC após análise da tratativa registrada pela
área responsável. \\
\hline
RN09 & Quando houver transferência de responsável, a
alteração deve preservar a rastreabilidade do chamado e
notificar o novo responsável. \\
\hline
RN10 & Em situações extraordinárias de impacto imediato,
como indisponibilidade generalizada de sistemas, o processo
deve permitir tratamento prioritário sem impedir o registro
formal da ocorrência. \\
\hline
\end{tabular}
\raggedright
\par
{\fontsize{10}{12}\selectfont Fonte: Elaborado pelo autor
com base na pesquisa de campo}
\end{table}

\subsection{Decisões de Modelagem Derivadas dos Requisitos}

Durante a transformação dos requisitos em solução, algumas
decisões de modelagem foram adotadas para reduzir erros de
direcionamento e separar responsabilidades. Embora a
Analista tenha indicado ``Área responsável'' entre as
informações desejáveis na abertura, a mesma respondente
também solicitou que a área de Clientes/Qualidade atuasse
como filtro interno para identificar o destino correto.
A Supervisora, por sua vez, apontou a necessidade de
transferência quando o direcionamento inicial estivesse
incorreto. Em razão dessa convergência, a definição da área
responsável, da prioridade, da tipologia e do responsável
pela tratativa foi posicionada na etapa de classificação,
sob responsabilidade da Qualidade/SAC, e não no formulário
inicial do solicitante.

Também foi estabelecida uma etapa intermediária entre a
tratativa e o encerramento. Após o preenchimento da causa
raiz e da ação corretiva, o chamado deve retornar para
análise da Qualidade/SAC. Essa decisão atende à necessidade
de obrigatoriedade de posicionamento, análise de causa e
solução indicada no levantamento e impede que a própria área
responsável encerre diretamente um chamado sem validação.

Por fim, o levantamento demonstrou que a automação não deve
ser compreendida apenas como mecanismo de registro. A
necessidade de governança, rastreabilidade e indicadores
foi recorrente nos três níveis consultados, tornando a
centralização dos dados e a manutenção de um histórico
estruturado elementos centrais da solução.

\section{Análise e Redesenho do Processo}

Esta seção materializa a etapa de \textit{Redesenho} do ciclo
de BPM apresentado na Base Teórica, articulando o diagnóstico
do processo atual (AS-IS), descrito na seção de Levantamento
de Requisitos, com a especificação do processo futuro
(TO-BE), que orienta o desenvolvimento da solução apresentado
na seção seguinte. O objetivo não é apenas contrastar duas
representações do fluxo, mas evidenciar, campo a campo,
decisão a decisão, a origem de cada mudança nos requisitos
levantados junto aos três perfis consultados.

\subsection{Processo Atual (AS-IS)}

O diagnóstico consolidado a partir das respostas da Analista,
da Supervisora e do Gerente caracteriza um processo
descentralizado e dependente de controles manuais. Os
chamados internos eram registrados no SAP Service --- mesma
ferramenta utilizada como CRM para atendimentos e reclamações
de clientes externos ---, misturando ocorrências internas e
manifestações externas em um único ambiente. A chegada do
chamado era comunicada paralelamente por um grupo de
WhatsApp, sem protocolo automático, sem controle de prazo e
sem visibilidade gerencial sobre volume, prioridade ou tempo
de resolução. O Quadro~\ref{quadro:as-is} sintetiza as
características centrais do processo atual.

\begin{table}[H]
\caption{Características do processo atual (AS-IS)}
\label{quadro:as-is}
\centering
\begin{tabular}{p{3.6cm} p{8.2cm}}
\hline
\textbf{Dimensão} & \textbf{Situação observada} \\
\hline
Canal de abertura & Cadastro ``Experiência do Cliente'' no
SAP Service, compartilhado com manifestações externas de
clientes. \\
\hline
Comunicação & Complementar e manual, por grupo de WhatsApp,
sem registro formal do repasse. \\
\hline
Controle de prazo & Inexistente; sem SLA definido ou
monitorado. \\
\hline
Rastreabilidade & Baixa; mistura entre chamados internos e
externos dificulta a leitura de volumes. \\
\hline
Responsabilidade & Informal; dependência de retorno manual
das áreas responsáveis. \\
\hline
Indicadores & Inexistentes; ausência de base para gestão
baseada em dados. \\
\hline
\end{tabular}
\raggedright
\par
{\fontsize{10}{12}\selectfont Fonte: Elaborado pelo autor com
base na pesquisa de campo}
\end{table}

\subsection{Processo Futuro (TO-BE)}

O processo redesenhado responde a cada uma das fragilidades
identificadas no AS-IS, substituindo controles manuais por
regras automatizadas e centralizando o registro em uma única
fonte de dados. O Quadro~\ref{quadro:to-be} apresenta a
transformação proposta, dimensão a dimensão, em correspondência
direta com o Quadro~\ref{quadro:as-is}.

\begin{table}[H]
\caption{Transformação do processo: AS-IS para TO-BE}
\label{quadro:to-be}
\centering
\begin{tabular}{p{3.0cm} p{4.9cm} p{4.9cm}}
\hline
\textbf{Dimensão} & \textbf{AS-IS} & \textbf{TO-BE} \\
\hline
Canal de abertura & Cadastro compartilhado no SAP Service &
Site no SharePoint com aplicativo Power Apps próprio para
chamados internos \\
\hline
Comunicação & Grupo de WhatsApp, sem registro & Notificações
automáticas por e-mail, vinculadas a eventos da Microsoft
List \\
\hline
Controle de prazo & Inexistente & SLA por tipo e prioridade,
com contagem automatizada \\
\hline
Rastreabilidade & Baixa, dados misturados & Protocolo único,
histórico centralizado na Microsoft List \\
\hline
Responsabilidade & Informal & Classificação e atribuição
explícita pela Qualidade/SAC, com transferência rastreável \\
\hline
Indicadores & Inexistentes & Painel gerencial no Power BI
(volume, tempo médio, cumprimento de SLA) \\
\hline
\end{tabular}
\raggedright
\par
{\fontsize{10}{12}\selectfont Fonte: Elaborado pelo autor}
\end{table}

\subsection{Decisões de Redesenho}

A transição entre os dois estados do processo não decorreu da
transposição direta do fluxo manual para um ambiente digital,
mas de decisões de redesenho fundamentadas nos requisitos
levantados. Duas decisões merecem destaque.

A primeira separa o registro de chamados internos das
manifestações externas de clientes, respondendo diretamente à
fragilidade central identificada pela Analista no processo
AS-IS. Essa separação é também uma decisão de Engenharia de
Processos: ao isolar o fluxo interno em um canal próprio, a
organização ganha capacidade de leitura de volumes e de
governança que o modelo anterior não permitia.

A segunda desloca a definição de área responsável, prioridade,
tipologia e responsável pela tratativa do formulário inicial
de abertura para uma etapa intermediária de classificação,
sob responsabilidade da Qualidade/SAC. Essa decisão resulta da
convergência entre a necessidade de filtro interno apontada
pela Analista e a necessidade de transferência de responsável
apontada pela Supervisora, e ilustra como o redesenho do
processo (TO-BE) é uma consequência direta do levantamento de
requisitos, e não uma escolha arbitrária de arquitetura de
software.

Por fim, foi estabelecida uma etapa de retorno obrigatório à
Qualidade/SAC após o preenchimento da causa raiz e da ação
corretiva pela área responsável, antes do encerramento
definitivo do chamado. Essa regra opera como mecanismo de
validação previsto pelo ciclo de BPM, impedindo que a própria
área executora encerre um chamado sem revisão, e reforça a
rastreabilidade e a governança como objetivos centrais do
redesenho.

O processo futuro (TO-BE) especificado nesta seção é a base
direta para os componentes desenvolvidos e descritos na seção
de Desenvolvimento da Solução, na qual cada estado, campo e
regra de negócio aqui definidos é implementado no ecossistema
Microsoft 365.

\section{Desenvolvimento da Solução}

Esta seção corresponde à etapa de \textit{Automação} do
ciclo de BPM apresentado na Base Teórica: o processo futuro
(TO-BE) especificado na seção de Análise e Redesenho do
Processo é aqui transformado em componentes executáveis do
ecossistema Microsoft 365, preservando, mesmo em ambiente
\textit{low-code}, as etapas de arquitetura, implementação e
validação próprias da Engenharia de Software. O
desenvolvimento foi realizado com base nos requisitos
consolidados na etapa anterior e utilizou componentes do
ecossistema Microsoft 365 já disponível na organização. A
solução resultante foi estruturada em cinco elementos
integrados: uma Central de Chamados no SharePoint, um
aplicativo Canvas no Power Apps para abertura das
solicitações, uma Microsoft List denominada
\textit{Chamados} como repositório central, dois fluxos
principais no Power Automate e uma área de acesso aos
indicadores gerenciais.

A arquitetura implementada difere parcialmente da
arquitetura inicialmente descrita na proposta. Na versão
desenvolvida, o SharePoint atua como portal de acesso,
enquanto o registro do chamado é realizado por um
aplicativo Power Apps incorporado à página. O Power Apps
grava diretamente o novo item na Microsoft List e, a partir
da criação desse item, o fluxo FL01 é disparado para gerar o
protocolo e realizar as notificações. Essa alteração
preservou o SharePoint como ponto único de entrada, porém
separou a camada de interface da camada de armazenamento,
permitindo maior controle sobre a experiência do
solicitante.

\subsection{Arquitetura da Solução Desenvolvida}

A Tabela~\ref{tab:arquitetura-desenvolvida} apresenta os
componentes efetivamente utilizados e suas responsabilidades
na versão desenvolvida.

\begin{table}[H]
\caption{Componentes da arquitetura desenvolvida}
\label{tab:arquitetura-desenvolvida}
\centering
\begin{tabular}{p{3.0cm} p{4.1cm} p{5.1cm}}
\hline
\textbf{Componente} & \textbf{Função principal} &
\textbf{Responsabilidade no fluxo} \\
\hline
SharePoint &
Portal corporativo da solução &
Centralizar o acesso para abertura de chamado, consulta de
registros e acesso ao painel gerencial. \\
\hline
Power Apps &
Interface de abertura &
Coletar os dados do solicitante, validar campos
obrigatórios, permitir anexos e gravar o novo chamado na
Microsoft List. \\
\hline
Microsoft List &
Repositório central &
Armazenar dados de abertura, classificação, responsável,
tratativa, status, datas e histórico de alterações. \\
\hline
Power Automate --- FL01 &
Automação da abertura &
Gerar protocolo, inicializar dados do chamado e notificar
Qualidade/SAC e solicitante. \\
\hline
Power Automate --- FL03 &
Automação da tratativa e encerramento &
Controlar notificações por mudança de status, transferência
de responsável, envio para análise, conclusão e
cancelamento. \\
\hline
Power BI &
Camada gerencial &
Disponibilizar a visualização dos indicadores e da base
histórica a partir do portal da Central de Chamados. \\
\hline
\end{tabular}
\raggedright
\par
{\fontsize{10}{12}\selectfont Fonte: Elaborado pelo autor}
\end{table}

O fluxo lógico da solução pode ser representado pela
sequência: solicitante acessa o SharePoint, realiza a
abertura no Power Apps, o registro é persistido na
Microsoft List, o FL01 gera protocolo e notificações, a
Qualidade/SAC classifica e direciona o chamado, a área
responsável realiza a tratativa e o FL03 controla as
mudanças subsequentes até o encerramento.

\subsection{Central de Chamados no SharePoint}

O SharePoint foi utilizado como camada de apresentação e
ponto central de acesso ao processo. A página desenvolvida
apresenta uma identificação institucional da Central de
Chamados Internos e disponibiliza acessos para abertura de
novo chamado, consulta dos chamados registrados e acesso ao
dashboard gerencial. Essa organização concentra em um único
ambiente as funções que, no processo anterior, estavam
distribuídas entre SAP Service, mensagens e comunicações
paralelas.

A utilização do SharePoint como portal também permite que o
processo seja incorporado ao ambiente corporativo já
utilizado pelos usuários, reduzindo a necessidade de
acesso a uma aplicação externa. O formulário de abertura
foi incorporado à página por meio de um aplicativo Power
Apps, enquanto a Microsoft List permanece responsável pelo
armazenamento e pelas etapas internas de classificação,
tratativa e encerramento.

\subsection{Aplicativo Power Apps para Abertura do Chamado}

A primeira etapa do formulário originalmente prevista
diretamente na Microsoft List foi substituída por um
aplicativo Canvas desenvolvido no Power Apps. A mudança foi
realizada com o objetivo de oferecer uma interface mais
simples ao solicitante e separar os campos de abertura dos
campos internos de classificação e tratamento.

O aplicativo foi integrado diretamente à lista
\textit{Chamados} e configurado para criação de novos
registros. Na versão analisada, o formulário de abertura
contém os campos Empresa, Tipo de chamado, Assunto do
chamado, Problema, Linha, Veículo e Anexos. Empresa, Tipo
de chamado, Assunto do chamado e Problema foram
configurados como obrigatórios. O campo Problema utiliza
entrada de texto em múltiplas linhas, possibilitando uma
descrição detalhada da ocorrência, enquanto o controle de
Anexos permite incluir evidências que apoiem a análise.

Os campos de escolha são carregados a partir das opções
cadastradas na própria Microsoft List, evitando a
manutenção de listas paralelas no aplicativo. O botão
\textit{Criar Chamado} utiliza a validação do formulário
antes do envio, enquanto o botão \textit{Limpar Formulário}
reinicializa os campos para uma nova entrada.

Após o salvamento bem-sucedido, o aplicativo altera seu
estado visual e apresenta ao usuário uma confirmação de que
o chamado foi aberto e está aguardando análise da
Qualidade/SAC. A mesma tela disponibiliza o botão
\textit{Novo Chamado}, que limpa o formulário, retorna ao
modo de criação e permite novo registro. Em caso de falha
no envio, uma mensagem de erro é apresentada ao usuário.

A interface foi ajustada para uso incorporado em página do
SharePoint, utilizando contêineres responsivos e dimensões
baseadas no espaço disponível. Essa adequação buscou
preservar a legibilidade do formulário quando executado
dentro do \textit{web part} do Power Apps.

\subsection{Estrutura da Microsoft List}

A Microsoft List denominada \textit{Chamados} constitui o
repositório central da solução. A estrutura separa os dados
informados pelo solicitante dos dados preenchidos durante o
processamento interno.

Na abertura são armazenados, entre outros, Empresa, Tipo de
chamado, Assunto, Problema, Veículo, Linha, anexos e o
usuário responsável pela criação do item. Nas etapas
posteriores são utilizados os campos Prioridade, Área
Responsável, Tipologia, Responsável pela tratativa, Causa
raiz, Ação corretiva, Data da Tratativa, Status do Chamado,
dados de transferência, Data de Fechamento e controle de
notificação de fechamento.

O registro centralizado permite que todas as etapas
utilizem o mesmo item como referência, evitando a criação
de bases paralelas e mantendo o histórico de alterações
associado ao protocolo. A estrutura também sustenta as
automações do Power Automate, que consultam alterações de
campos específicos para identificar quando uma nova ação
deve ser executada.

\subsection{Fluxo FL01 --- Abertura de Chamados Internos}

O fluxo FL01 é iniciado automaticamente quando um novo item
é criado na Microsoft List. Sua primeira ação calcula o
protocolo utilizando o ano corrente e o identificador
numérico do item. O padrão implementado é
\texttt{C-AAAA-NNNN}. Assim, um item de identificador 14
criado em 2026, por exemplo, recebe o protocolo
\texttt{C-2026-0014}.

Após a geração do protocolo, o próprio item é atualizado.
Na versão exportada do fluxo, o status inicial é definido
como \textit{Aberto} e a prioridade é inicialmente
registrada como \textit{Média}. Em seguida são executadas
duas notificações distintas.

A primeira notificação é destinada à equipe de
Qualidade/SAC e informa que um novo chamado aguarda
classificação. O e-mail apresenta protocolo, solicitante,
empresa, tipo, assunto, linha, veículo, status e descrição
do problema, além de um link direto para o item. O objetivo
é iniciar a etapa de classificação e direcionamento.

A segunda notificação é destinada ao próprio solicitante e
funciona como confirmação formal de abertura. O e-mail
informa o número do protocolo, resume os dados registrados,
indica que a solicitação aguarda análise da Qualidade/SAC e
fornece o link para acompanhamento do chamado. Com isso, a
confirmação de abertura deixa de depender de comunicação
manual.

\subsection{Classificação e Direcionamento pela
Qualidade/SAC}

Após a abertura, a Qualidade/SAC realiza a classificação
diretamente no registro da Microsoft List. Nessa etapa são
definidos os campos Prioridade, Área Responsável, Tipologia
e Responsável pela tratativa. O status do chamado também é
controlado pela Qualidade/SAC.

A decisão de manter esses campos fora do formulário do
solicitante reduz o risco de encaminhamento incorreto e
formaliza a Qualidade/SAC como instância de triagem. Quando
o chamado é validado e direcionado para tratamento, o
status é alterado para \textit{Em tratativa}, iniciando a
execução das regras do FL03 associadas à atribuição do
responsável.

A estrutura da lista também contempla o campo Motivo do
Cancelamento. Quando a Qualidade/SAC decide não prosseguir
com a solicitação e altera o status para
\textit{Cancelado}, o motivo registrado passa a compor a
rastreabilidade do encerramento sem tratativa.

\subsection{Fluxo FL03 --- Tratativa e Encerramento}

O FL03 é disparado quando um item da lista é modificado.
Inicialmente, o fluxo obtém quais propriedades foram
alteradas e recupera o estado atual do item. Em seguida,
uma estrutura de decisão do tipo \textit{Switch} avalia o
status do chamado e direciona o processamento para o ramo
correspondente.

No estado \textit{Em tratativa}, o fluxo verifica se houve
mudança do status para identificar uma nova atribuição.
Quando essa condição é atendida, o responsável pela
tratativa recebe uma notificação por e-mail contendo os
dados necessários e o link para acesso ao chamado.

Ainda no estado \textit{Em tratativa}, o fluxo verifica se
o campo de responsável foi alterado. Quando ocorre uma
transferência, o registro é atualizado com informações de
controle da transferência e o novo responsável é
notificado. Essa regra foi incorporada após a necessidade
identificada no levantamento de requisitos de redistribuir
chamados sem perda de rastreabilidade.

Quando Causa raiz e Ação corretiva são preenchidas e as
respectivas propriedades são identificadas como alteradas,
o FL03 atualiza automaticamente o status para
\textit{Em Análise} e envia uma notificação à
Qualidade/SAC. A área responsável, portanto, registra a
tratativa, mas não realiza o encerramento definitivo.

A análise final é responsabilidade da Qualidade/SAC. Após
avaliar a tratativa, a equipe pode alterar manualmente o
status para \textit{Concluído} ou, quando aplicável,
\textit{Cancelado}. No caso de conclusão, o FL03 notifica o
solicitante e os responsáveis envolvidos e registra a data
de fechamento.

Para impedir que a própria atualização realizada pelo
fluxo gere sucessivas notificações de encerramento, foi
utilizado um campo de controle denominado
``Notificação de fechamento enviada?''. O ramo de
conclusão somente envia o e-mail quando esse campo indica
que a notificação ainda não foi realizada; na sequência, o
campo é atualizado para \textit{Sim}. Esse mecanismo atua
como proteção contra reprocessamento e reduz o risco de
\textit{loop} provocado pelo gatilho de modificação da
lista.

No ramo \textit{Cancelado}, o solicitante recebe uma
notificação específica informando que o chamado foi
cancelado pela Qualidade/SAC. A comunicação apresenta os
dados do chamado e o motivo do cancelamento registrado no
item, mantendo o histórico disponível para consulta.

\subsection{Estados do Chamado e Responsabilidades}

A implementação consolidou cinco estados principais:
\textit{Aberto}, \textit{Em tratativa}, \textit{Em Análise},
\textit{Concluído} e \textit{Cancelado}. A
Tabela~\ref{tab:estados-chamado} apresenta a função de cada
estado.

\begin{table}[H]
\caption{Estados implementados no fluxo de chamados}
\label{tab:estados-chamado}
\centering
\begin{tabular}{p{2.6cm} p{4.2cm} p{5.5cm}}
\hline
\textbf{Status} & \textbf{Responsável predominante} &
\textbf{Significado no processo} \\
\hline
Aberto &
Qualidade/SAC &
Chamado registrado e aguardando classificação e
direcionamento. \\
\hline
Em tratativa &
Área responsável &
Chamado classificado, atribuído e disponível para registro
da causa raiz e da ação corretiva. \\
\hline
Em Análise &
Qualidade/SAC &
Tratativa preenchida pela área responsável e aguardando
validação antes do encerramento. \\
\hline
Concluído &
Qualidade/SAC &
Tratativa validada e chamado encerrado, com notificação aos
envolvidos. \\
\hline
Cancelado &
Qualidade/SAC &
Chamado encerrado sem continuidade da tratativa, com motivo
registrado e comunicação ao solicitante. \\
\hline
\end{tabular}
\raggedright
\par
{\fontsize{10}{12}\selectfont Fonte: Elaborado pelo autor}
\end{table}

\subsection{Rastreabilidade e Notificações}

A rastreabilidade foi implementada por meio da combinação
entre protocolo único, item persistido na Microsoft List,
histórico de modificações e notificações vinculadas às
mudanças de estado. Cada chamado permanece associado ao
mesmo registro durante todo o ciclo, permitindo recuperar
os dados de abertura, classificação, responsável, tratativa
e encerramento.

As notificações foram utilizadas como mecanismo de
transição entre os responsáveis pelo processo. Na abertura,
Qualidade/SAC e solicitante são comunicados. Na atribuição
ou transferência, o responsável pela tratativa é
notificado. Após o registro da causa raiz e da ação
corretiva, a Qualidade/SAC recebe a solicitação para
análise. Na conclusão, solicitante e responsáveis são
informados do encerramento; no cancelamento, o solicitante
recebe a justificativa correspondente.

Esse desenho substitui a comunicação paralela utilizada no
processo anterior por mensagens vinculadas a eventos
registrados na própria base, reduzindo a dependência de
grupos de mensagens para informar movimentações do chamado.

\subsection{Aderência da Implementação aos Requisitos}

A análise entre requisitos e componentes desenvolvidos
demonstrou que a versão atual atende aos elementos centrais
de centralização, protocolo, classificação, atribuição,
transferência, tratativa, validação e encerramento.
Entretanto, alguns requisitos identificados na pesquisa de
campo não puderam ser confirmados como implementados nos
artefatos analisados e devem ser tratados como itens de
evolução ou de validação antes da conclusão definitiva do
projeto.

\begin{table}[H]
\caption{Aderência dos principais requisitos à versão desenvolvida}
\label{tab:aderencia-requisitos}
\centering
\begin{tabular}{p{1.2cm} p{6.6cm} p{4.4cm}}
\hline
\textbf{ID} & \textbf{Requisito} &
\textbf{Situação observada} \\
\hline
RF01--RF14 &
Centralização, abertura, protocolo, notificações,
classificação, atribuição, transferência, tratativa,
análise e encerramento. &
Implementados nos componentes SharePoint, Power Apps,
Microsoft List, FL01 e FL03. \\
\hline
RF15 &
Controle automático de SLA por tipo e prioridade. &
Requisito identificado na pesquisa, porém o cálculo
automático não foi evidenciado nos fluxos exportados
analisados. \\
\hline
RF16 &
Comentários, atualizações e justificativas durante
tratativas prolongadas. &
Necessidade registrada pela Supervisora; não foi
identificado campo específico na estrutura final analisada. \\
\hline
RF17 &
Quantidade de protocolos relacionados. &
Requisito informado pela Analista; não foi identificado no
formulário final de abertura analisado. \\
\hline
RF18 &
Tipologias reutilizáveis. &
Implementado na etapa de classificação da Microsoft List. \\
\hline
RF19 &
Indicadores gerenciais. &
A Central de Chamados possui acesso ao dashboard; a
validação detalhada das métricas deve ser realizada na
etapa de avaliação dos resultados. \\
\hline
\end{tabular}
\raggedright
\par
{\fontsize{10}{12}\selectfont Fonte: Elaborado pelo autor}
\end{table}

Outro ponto de aderência parcial está na regra de
priorização. A pesquisa de campo indicou prioridade Alta
para chamados Massivos e Média para chamados Internos.
Entretanto, na versão exportada do FL01, todo novo chamado
é inicialmente atualizado com prioridade Média, ficando a
classificação definitiva sob responsabilidade da
Qualidade/SAC. Dessa forma, a regra de prioridade por tipo
não está automatizada no momento analisado.

Esses itens não invalidam o fluxo principal desenvolvido,
mas representam diferenças entre os requisitos levantados
e a versão implementada. O registro explícito dessas
diferenças é importante para a avaliação acadêmica da
solução, pois permite distinguir funcionalidades
efetivamente entregues de requisitos planejados ou
reservados para evolução.
